\section{Conclusion}
\label{sec:conclusion}
A rank-and-weight index rule is recoverable from its published description. A
constrained-optimisation rule is recoverable only up to the parts of its
description that are written in a vocabulary the reader does not have. Our
replica of the heuristic index tracks its fund at 1.62\% annualised tracking
error; our replica of the optimised index tracks at 5.83\%, and 2.10 percentage
points of that gap turn on how one noun phrase in the constraint document is
read.

The usual explanation for the difference does not survive measurement. Four
covariance estimators, from a singular sample matrix to a structural factor
model, track the fund within 26 basis points of one another under one
specification, against 2.10 percentage points for the exemption reading; a
block bootstrap separates the two intervals cleanly and the ordering holds in
every resample. That is what prior work reports for constrained portfolios
generally, and the reason is also known: binding constraints act as shrinkage,
so a tightly constrained feasible set leaves estimator differences nowhere to
express themselves. Our contribution on that point is confirmation in a
specific published constraint set, not discovery. What is new here is the
comparison of magnitudes, and the observation that the larger of the two
quantities is a reading of a document rather than a property of an estimator.

What the constraints deliver is not lower realised risk, which over this sample
they conspicuously fail to deliver, but control over where risk accumulates.
The sector band binds at every review examined; without it, the heuristic
accumulates an underweight to one industry that grows monotonically past forty
per cent. Both designs are concentrated, in different dimensions, and only one
of them concentrates by stated rule. Neither index shows a six-factor alpha
significant at the five per cent level, and betting-against-beta loadings
between 8.7 and 14.0 absorb what intercept there is, which is consistent with
the leverage-constraint account and does not distinguish it from the
benchmarking one.

The practical implication is narrow enough to state exactly. For anyone
replicating a constrained index from public data, effort spent resolving what
the methodology document actually specifies is worth more than effort spent
improving the covariance estimate, by a factor of about eight in this case. The broader implication, which this evidence motivates rather than
establishes, is that an index methodology stated in terms of a proprietary
factor model is not fully reproducible by construction, and that the gap this
creates is measurable and worth measuring.
